\chapter{Trabalho a Ser Desenvolvido}
\label{cap:trabalho}

% Nesta seção, uma descrição funcional da solução proposta deverá ser feita:
% - O projeto deverá ser decomposto em módulos funcionais (incluindo software e hardware).
% - Deverá ser apresentado um diagrama em blocos mostrando uma visão geral do sistema a ser desenvolvido, isto é, como os módulos funcionais se relacionam.
% - Cada módulo funcional deve ser descrito individualmente.
% - Deve ser indicado também, de maneira sucinta, como cada módulo será implementado (linguagem de programação, hardware, etc.).

\section{Descrição Funcional e Diagrama em Blocos}

O objetivo do trabalho é desenvolver um software capaz de analisar
trechos de código-fonte antes de sua integração a um sistema,
identificando padrões associados a vulnerabilidades de segurança
conhecidas e classificando-os de forma contextual, de modo a
complementar as ferramentas tradicionais de análise estática e a
reduzir a taxa de falsos positivos por elas produzida.

Funcionalmente, a solução recebe como entrada um conjunto de trechos
de código-fonte -- tipicamente os arquivos modificados em um
\textit{pull request} ou em um \textit{commit} -- e devolve ao
desenvolvedor um relatório contendo os trechos considerados
potencialmente vulneráveis, o tipo de vulnerabilidade associado a cada
um e um nível de confiança da classificação, permitindo que os alertas
mais críticos sejam priorizados durante a revisão de código.

Para alcançar esse comportamento, o sistema foi decomposto em cinco
módulos funcionais, todos de \textit{software}, organizados como um
\textit{pipeline}: o módulo de entrada e pré-processamento (M1)
prepara o código recebido; o módulo de análise estática (M2) executa
as ferramentas tradicionais sobre esse código, gerando alertas brutos;
o módulo de extração de características (M3) converte o código
normalizado e os alertas em representações numéricas; o módulo de
classificação (M4) decide, a partir dessas representações, quais
trechos são de fato vulneráveis e qual a criticidade de cada alerta; e
o módulo de relatório e interface (M5) consolida e apresenta os
resultados ao desenvolvedor. O projeto não contempla módulos de
\textit{hardware} dedicados: a execução ocorre integralmente em
computadores de propósito geral, conforme detalhado na
Seção~\ref{sec:hardware}.

A Figura~\ref{fig:blocos} apresenta a visão geral do sistema e o
relacionamento entre esses módulos. As caixas em traço contínuo
correspondem aos módulos a serem desenvolvidos, as caixas em traço
tracejado representam recursos externos consumidos pelo sistema
(ferramentas de análise estática de terceiros e o modelo treinado
junto ao conjunto de dados de vulnerabilidades), e as caixas em
destaque representam as fronteiras do sistema, isto é, a origem do
código analisado e o usuário final da solução.

\begin{figure}[htbp]
\centering
\begin{tikzpicture}[
    node distance=1.4cm and 1.9cm,
    mod/.style={draw, rounded corners, minimum width=3.2cm, minimum height=1.25cm, align=center, font=\small, fill=blue!12},
    ext/.style={draw, dashed, rounded corners, minimum width=3.1cm, minimum height=1.05cm, align=center, font=\footnotesize, fill=black!6},
    io/.style={draw, rounded corners=2pt, minimum width=2.9cm, minimum height=1.25cm, align=center, font=\small, fill=green!12},
    arr/.style={-{Latex[length=2.2mm]}, thick},
    lbl/.style={font=\scriptsize, midway, align=center}
]
    % ---------- Linha 1 ----------
    \node[io] (src) {Código-Fonte\\\footnotesize(repositório / PR)};
    \node[mod, right=of src] (m1) {\textbf{M1}\\Entrada e\\Pré-processamento};
    \node[mod, right=of m1] (m2) {\textbf{M2}\\Análise Estática};
    \node[ext, above=0.8cm of m2] (tools) {SonarQube, linters,\\Dependabot};

    % ---------- Linha 2 ----------
    \node[mod, below=1.9cm of m2] (m3) {\textbf{M3}\\Extração de\\Características};
    \node[mod, left=of m3] (m4) {\textbf{M4}\\Classificação\\(ML / LLM)};
    \node[mod, left=of m4] (m5) {\textbf{M5}\\Relatório e\\Interface};
    \node[ext, below=0.9cm of m4] (model) {Modelo treinado\\+ dataset CWE/CVE};
    \node[io, below=0.9cm of m5] (dev) {Desenvolvedor};

    % ---------- Fluxo principal ----------
    \draw[arr] (src) -- (m1);
    \draw[arr] (m1) -- node[lbl, above]{código\\normalizado} (m2);
    \draw[arr] (m2) -- node[lbl, right, xshift=1mm]{alertas\\brutos} (m3);
    \draw[arr] (m3) -- node[lbl, above]{\textit{features}} (m4);
    \draw[arr] (m4) -- node[lbl, above]{alertas\\classificados} (m5);
    \draw[arr] (m5) -- (dev);

    % ---------- Entradas auxiliares ----------
    \draw[arr, dashed] (tools) -- (m2);
    \draw[arr, dashed] (model) -- (m4);

    % ---------- Código normalizado também alimenta M3 ----------
    \draw[arr] (m1.south) -- ++(0,-0.95) -| node[lbl, below, pos=0.25]{código normalizado} (m3.north);
\end{tikzpicture}
\caption{Diagrama em blocos do sistema.}
\label{fig:blocos}
\end{figure}

O fluxo de dados entre os módulos ocorre da seguinte forma. O código
recebido é normalizado por M1, que o entrega tanto a M2 quanto a M3 --
este último precisa do código em si, e não apenas dos alertas, para
extrair características estruturais. M2 executa as ferramentas de
análise estática e repassa a M3 os alertas brutos resultantes. M3
combina as duas fontes em um vetor de características por trecho de
código, que é consumido por M4. Este, apoiado no modelo previamente
treinado, produz para cada trecho uma classificação e um grau de
confiança, os quais são finalmente formatados por M5 e apresentados ao
desenvolvedor.

\section{Detalhamento dos Módulos Funcionais}

\subsection{M1 -- Entrada e Pré-processamento}

Responsável por receber o código-fonte a ser analisado e prepará-lo
para as etapas seguintes. O módulo obtém os arquivos a partir de um
repositório Git (ou dos arquivos alterados em um \textit{pull
request}), identifica a linguagem de cada arquivo, normaliza aspectos
superficiais do texto -- codificação, quebras de linha e indentação --
e realiza o \textit{parsing} do código, produzindo uma representação
estruturada que permite localizar funções, chamadas, variáveis e
dependências. Cabe também a este módulo delimitar o escopo da análise,
selecionando apenas os trechos novos ou modificados, já que o objetivo
da solução é avaliar o código antes de sua integração.

\textit{Implementação:} Python, utilizando bibliotecas de manipulação
de repositórios Git e de \textit{parsing} de código-fonte para a
construção das árvores sintáticas. A saída é entregue aos módulos
seguintes em formato estruturado (JSON).

\subsection{M2 -- Análise Estática}

O módulo orquestra a execução de ferramentas de Análise estática, coleta os alertas por elas produzidos e os converte para um formato unificado, de modo que alertas de
ferramentas distintas possam ser tratados de maneira homogênea pelos
módulos seguintes. Cada alerta preserva a localização no código
(arquivo, linha e função), a categoria da vulnerabilidade e a
severidade atribuída pela ferramenta de origem.

\textit{Implementação:} Python para a orquestração e a normalização
dos resultados, integrando-se às ferramentas SonarQube, a
\textit{linters} específicos de cada linguagem e ao Dependabot. As
ferramentas externas são executadas em contêineres Docker, o que
garante a reprodutibilidade do ambiente de análise.

\subsection{M3 -- Extração de Características}

Responsável por converter as entradas textuais e estruturais em
representações numéricas processáveis por algoritmos de aprendizado de
máquina. O módulo combina duas fontes: o código normalizado vindo de
M1, do qual são extraídas características estruturais e semânticas do
trecho analisado, e os alertas vindos de M2, dos quais são derivadas
características relativas ao tipo, à severidade e à quantidade de
alertas incidentes sobre aquele trecho. O resultado é um vetor de
características por trecho de código, no qual cada alerta candidato
está associado ao seu respectivo contexto.

\textit{Implementação:} Python, com bibliotecas de Machine Learning
para vetorização e, alternativamente, representações baseadas em
\textit{embeddings} de modelos pré-treinados de código, conforme
avaliação experimental a ser realizada.

\subsection{M4 -- Módulo de Classificação (ML / LLM)}

Núcleo da solução, responsável por decidir quais trechos apresentam de
fato risco de vulnerabilidade e qual a criticidade de cada alerta. A
partir dos vetores de características recebidos de M3, o modelo de
classificação supervisionada, treinado sobre um conjunto
de dados de vulnerabilidades conhecidas; Complementarmente, Large Language Models
poderão ser empregados como apoio à análise contextual dos casos de
maior ambiguidade, avaliando a intenção do trecho de código e as
relações entre funções que escapam às regras sintáticas das
ferramentas tradicionais. Trata-se, portanto, do módulo que efetiva o
diferencial do projeto em relação à análise estática convencional.

\textit{Implementação:} Python, com bibliotecas de Machine Learning
para o treinamento e a inferência dos modelos de classificação; o
apoio de Large Language Models será realizado por meio de consultas a
modelos de linguagem, tratado como componente complementar e opcional
do \textit{pipeline}.

\subsection{M5 -- Relatório e Interface}

Responsável por consolidar os resultados da análise e apresentá-los ao
desenvolvedor de forma compreensível e acionável. O módulo agrega os
alertas classificados por M4, ordena-os por criticidade e gera um
relatório indicando, para cada ocorrência, o arquivo e a linha
correspondentes, o trecho de código envolvido, o tipo de
vulnerabilidade sugerido e o nível de confiança atribuído pelo modelo.
A interface é o ponto de contato entre a solução e seu público-alvo,
razão pela qual se prioriza a clareza da informação sobre a
exaustividade dos dados exibidos.

\textit{Implementação:} Python, com geração de relatório, de modo
a permitir a execução tanto manual quanto integrada a um
\textit{pipeline} de integração contínua.

\subsection{Considerações sobre \textit{Hardware}}
\label{sec:hardware}

O projeto não prevê o desenvolvimento de módulos de \textit{hardware}.
A execução da solução ocorre em computadores de propósito geral: os
computadores pessoais dos integrantes serão utilizados para o
desenvolvimento e para os testes do \textit{pipeline}, e
infraestrutura computacional local ou em nuvem poderá ser empregada
quando o treinamento dos modelos de Machine Learning exigir maior
capacidade de processamento, em especial na presença de GPU.
\end{document}

